This problem can be solved by using relationship between the telescope's
resolution and the angular size of black holes which can be seen in this
equation:
\[
\phi_{\mathrm{telescope}} \leq \phi_{\mathrm{blackhole}}
\quad\Rightarrow\quad
1.2\frac{\lambda}{D_{\mathrm{telescope}}}
 \leq \frac{\text{Diameter of event horizon}}
           {\text{Observer distance to the Galactic center}}
\]
\[
1.2\frac{\lambda}{D_{\mathrm{telescope}}}
 \leq \frac{2R_{\mathrm{blackhole}}}{d}
\]
\hfill(25)

Then, it is assumed that the radius of black hole is the Schwarzschild
radius; and the telescope is also Earth-sized telescope. Therefore:
\[
1.2\frac{\lambda}{2R_{\mathrm{Earth}}}
 \leq \frac{2R_{\mathrm{Schwarzschild,BH}}}{d}
\]
\begin{align*}
d &\leq \frac{4 \times R_{\mathrm{Schwarzschild,BH}} \times R_{\mathrm{Earth}}}
             {1.2 \times \lambda}
\quad\Rightarrow\quad
d \leq \frac{4 \times M_{\mathrm{Black\,Hole}}
             \times R_{\mathrm{Schwarzschild,Sun}} \times R_{\mathrm{Earth}}}
            {1.2 \times M_{\mathrm{Sun}} \times \lambda}\\
d &\leq \frac{4 \times 2.1 \times 10^{10}\,M_{\mathrm{Sun}}
             \times R_{\mathrm{Schwarzschild,Sun}} \times R_{\mathrm{Earth}}}
            {1.2 \times M_{\mathrm{Sun}} \times \lambda}
\quad\Rightarrow\quad
d \leq \frac{8.4 \times 10^{10}
             \times R_{\mathrm{Schwarzschild,Sun}} \times R_{\mathrm{Earth}}}
            {1.2 \times \lambda}
\end{align*}
\hfill(50)

To find the maximum distance, then we should use the minimum $\lambda$.
Therefore:
\[
d_{\max} = \frac{8.4 \times 10^{10}
            \times R_{\mathrm{Schwarzschild,Sun}} \times R_{\mathrm{Earth}}}
           {1.2 \times \lambda_{\min}}
 = \frac{8.4 \times 10^{10} \times 2.95 \times 10^{3}
         \times 6.3708 \times 10^{6}}
        {1.2 \times 20 \times 10^{-6}}\ \mathrm{m}
 = 6.6 \times 10^{25}\ \mathrm{m}
\]
\hfill(25)
